\section{Discussion}
\label{sec:discussion}

\ourbench{} separates two empirical questions: the recoverability of cell
identity at different hierarchy levels, and the incremental contribution of
paired APT beyond a specified RNA reference. The APT-only oracle analyses now
separate three explanations for the large ordinary-oracle gap. Training-only
lineage priors do not explain it; neither does retraining after preserving
patient-by-lineage APT distributions while breaking cell-level correspondence.
However, a simple trained conditional subtype MLP does not yield a resolved
overall deployable gain. A targeted follow-up that directly penalizes fine-head
parent mass, aligns coarse and fine-implied lineage distributions, and distills
a true-lineage teacher reaches at most a small unresolved point gain; it does
not reduce cross-lineage errors. The paired-input increment asks a different question:
it persists over the two tested RNA widths, but the RNA+APT advantage over the
repeated patient-by-lineage-shuffle control remains unresolved overall.
Improvements over RNA alone therefore do not establish finer-grained pairing
specificity. Patient-disjoint testing and subject-balanced scoring provide the
common evaluation foundation, not an additional claim of methodological novelty.

\paragraph{A reusable comparison setting.}
The resource specifies an APT input, operational ontology, fixed patient folds,
metrics, and reference results. Soft-cascade is a hierarchy-aware reference,
not a claimed algorithmic advance. These references establish comparison
points for improving APT cell typing without changing the input modality,
ontology, or held-out patient partitions. Complete data access, permissions,
annotation provenance, and a portable evaluator remain necessary release
requirements. The unified comparison supplies full-cell-budget classical and neural
references, with separate structure/supervision contrasts. Neither the
Flat-minus-MLP nor Cascade-minus-Flat interval establishes superiority or
equivalence. More comprehensive APT-compatible annotation baselines remain
useful extensions.

\paragraph{Boundaries of the evidence.}
The APT cohort contains \numparticipants{} participants and RNA-derived reference labels.
Independent annotation checks and additional APT cohorts are needed to
assess broader generalization. Patient-disjoint testing does not remove
disease-correlated acquisition effects; this manuscript makes no disease
classification or clinical-validation claim. RNA supplies the reference labels, and its apparent
difficulty depends on features, preprocessing, and the fitted reference.
An APT increment over either 2{,}000 or 5{,}000 HVGs may reflect restricted RNA representation;
it is not evidence of irreducible information absent from RNA. The capped
probe does not establish an increment over a full-cell-budget,
capacity-controlled RNA reference. Likewise, oracle-versus-prior performance
does not guarantee that predicted-lineage decoding realizes the oracle gain.
The v2 frozen-score tests address two deployable readout rules; the conditional
subtype MLP and HBM objectives test a limited set of trained alternatives,
without covering all possible hierarchical learners. True-lineage masking and shuffling are
privileged diagnostics, and the latter does not identify a causal or molecular
mechanism. Three training seeds and
pointwise resampling intervals also provide limited uncertainty evidence.
The outer folds have been examined repeatedly during this study; fixed
follow-up analyses on them are not untouched confirmatory tests. Their
overlapping development sets also limit interpreting OOF bootstrap intervals
as uncertainty for a learning algorithm on an independently sampled cohort.

\section{Conclusion}
\label{sec:conclusion}

\ourbench{} establishes a patient-disjoint evaluation setting for cell typing
from single-cell aptamer profiles, with a fixed label ontology, explicit
subject-balanced scoring, and classical and neural reference results.
It asks which hierarchy levels are recoverable in new patients and whether
APT contributes additional typing information beyond an RNA reference.
Training-prior-controlled oracle decoding reveals usable but heterogeneous
within-lineage discrimination under the tested APT-only models, and the
patient-by-lineage shuffle control supports cell-level APT correspondence in
that privileged diagnostic setting. The tested conditional subtype MLP does
not translate this signal into an overall deployable gain, nor do direct
fine-head parent supervision, coarse--fine consistency, or privileged
within-lineage distillation resolve the gap. The repeated
paired RNA+APT probe improves over the tested RNA-only references at both
widths, while its overall advantage over patient-by-lineage-shuffled APT
remains unresolved. Distinguishing these findings avoids equating improved
reference reconstruction with established subtype-specific complementarity.
